\def\UseOption{print,lngcs}
\expandafter\ifx\csname classoptions\endcsname\relax
  \def\classoptions{}
\fi
\documentclass[fykos,\classoptions]{fksgeneric}% TODO hlavička, asi něco jako mfnaboj
\usepackage{pst-barcode}

%\pagestyle{empty}
\renewcommand\makefooter{}

\fancyhead[L]{\bfseries\small FYKOSí Fyziklání ROK}
\fancyhead[R]{}
\fancyfoot{}


\setcounter{year}{3}%tady bylo TODO
\setcounter{batch}{1}

%%%%%%%%%%%%%%%%%%%%%%%%%%%%%%%%%%%%%%%%%%%%%%%%%%%%%%%%%%%
%%%%%%%%%%%%%%%%%%%%%%% Makra %%%%%%%%%%%%%%%%%%%%%%%%%%%%%
%%%%%%%%%%%%%%%%%%%%%%%%%%%%%%%%%%%%%%%%%%%%%%%%%%%%%%%%%%%
\def\city{}
\newcount\team\team=\teamscount \multiply\team by 5
\newcount\teamN\teamN=0%
%\newcount\pps\pps=20% pages per 4survey; 4str, 5 clenu

\newcounter{boxcnt}
\newcounter{sboxcnt}
\newcounter{questioNum}
\setcounter{questioNum}{0}

% velikosti boxu na odpoved
\newlength{\AboxH}\setlength{\AboxH}{1.5\baselineskip}
\newlength{\AboxD}\setlength{\AboxD}{0.5\baselineskip}

% pocet ctverecku
\newcounter{sqCnt}

\newcommand{\nl}[0]{&\ifnum\theboxcnt>0\addtocounter{boxcnt}{-1}%
\setcounter{sboxcnt}{\thesqCnt}%
\loop\ifnum\thesboxcnt>0\addtocounter{sboxcnt}{-1}%
\fbox{\vrule width 3pt height 0pt depth 0pt \vrule width 0pt height 3pt depth 0pt}\ifnum\thesboxcnt>0\hfill\fi\repeat%
\else \fbox{\vrule width 3.53cm height 0pt depth 0pt %3.53cm je tak akorat, aby se vesla v pravemu okraji
\vrule width 0pt height \AboxH depth \AboxD}\fi\smallskip\\}

% otazka; odpovedi oddělení \nl; počet checkboxů (další jsou boxy)
\newcommand{\question}[4]{%
\refstepcounter{questioNum}\addtocounter{questioNum}{-1}%
\setcounter{sqCnt}{1}%
\noindent\begin{minipage}[t]{8mm}\stepcounter{questioNum}\bf\arabic{questioNum}.\end{minipage}%
\begin{minipage}[t]{3cm}{\sl #1}\end{minipage}\hspace{2mm}%
\setcounter{boxcnt}{#3}%
\begin{minipage}[t]{\textwidth-8mm-3cm-2mm-0.3pt}%8mm číslo, 3cm otázka, 2mm skvíra, .3pt je z Bulharska
\begin{tabular}[t]{p{5cm}l}
#2\nl%
\end{tabular}\end{minipage}%\vrule%
\par#4\par\nopagebreak\vfil\smallskip\hrule\vfil\vfil\smallskip%
}

\newcommand{\evaluation}[4]{%
\refstepcounter{questioNum}\addtocounter{questioNum}{-1}%
\setcounter{sqCnt}{5}%
\noindent\begin{minipage}[t]{8mm}\stepcounter{questioNum}\bf\arabic{questioNum}.\end{minipage}%
\begin{minipage}[t]{3cm}\sl #1\end{minipage}\hspace{2mm}%
\setcounter{boxcnt}{#3}%
\raisebox{\baselineskip}{% posun nahoru o nadpis
\begin{minipage}[t]{\textwidth-8mm-3cm-2mm-0.3pt}\begin{tabular}[t]{p{5cm}l}%
& \emph{\hspace{1.5pt}1\hfill2\hfill3\hfill4\hfill5\hspace{1.5pt}~}\\%
#2\nl%
&\vrule width 3.78cm height 0pt depth 0pt
\end{tabular}\end{minipage}}%
\vspace{-\baselineskip}
\par#4\par\nopagebreak\vfil\smallskip\hrule\vfil\vfil\smallskip%
}

\newcommand{\goto}[1]{\emph{\scriptsize\opt{lngcs}{Pokračujte na otázku}\opt{lngsk}{Pokračujte na otázku} \ref{#1}.}}

%%%%%%%%%%%%%%%%%%%%%%%%%%%%%%%%%%%%%%%%%%%%%%%%%%%%%%%%%%%
%%%%%%%%%%%%%%%%%%%% Zacatek ankety %%%%%%%%%%%%%%%%%%%%%%%
%%%%%%%%%%%%%%%%%%%%%%%%%%%%%%%%%%%%%%%%%%%%%%%%%%%%%%%%%%%
\begin{document}

\loop\ifnum\teamN<\the\team%
\setcounter{questioNum}{0}%
\section{\opt{lngcs}{Anketa pro účastníky FYKOSího Fyziklání ROK}}
Otázky vyplňujte začerněním příslušného políčka. Pokud 
si přejete změnit odpověď, začerněné políčko škrtněte velkým křížkem 
a~začerněte jiné. Otázky prosíme vyplňujte popořadě a~respektujte skoky. {\bf Anketu nepřehýbejte!}%


%opal MFF
%TODO musí odpovídat na všechno?
%TODO mohou použít modrou tužku; na MFF jsem to udělal.

\opt{lngcs}{Vzor pro označování:}
\begin{center}
\opt{lngcs}{neoznačeno}\quad \includegraphics{neozn.eps}\quad% 
\opt{lngcs}{označeno}\quad \includegraphics{ozn.eps}\quad%
\opt{lngcs}{opraveno}\quad \raisebox{-3pt}{\includegraphics{opr.eps}}\quad%
\end{center}


%Pro vyplnění ankety použijte modrou či černou propisovací tužku.
%Pokud chtete označit nějakou odpověď, vybarvěte pečlivě 
%odpovídající čtvereček. Pokud chcete Vaše rozhodnutí změnit, nejprve 
%přešktrněte původní volbu a~pak zabarvěte čtvereček odpovídající nové
%volbě.

%Manuál viz naše anketa. Vyplnit kolečko, pak jej přeškrtnout.\medskip
%Nebolo by lepsie tam dat vzor ako na MFF ankete? Su to deti a to od MFF je nazornejsie. //Paťo
%Bylo, ale nechtělo se mi/neměl jsem čas kreslit obrázky (udělám to) //Lukáš

\subsection{\opt{lngcs}{Korespondenční semináře}}
\question{\opt{lngcs}{Jsem studentem ročníku:}}{%
    \opt{lngcs}{1.}\nl 
    \opt{lngcs}{2.}\nl 
    \opt{lngcs}{3.}\nl 
    \opt{lngcs}{4.}}{4}{}
\question{\opt{lngcs}{Fyziky a~matematiky mám ve škole:}}{%
    \opt{lngcs}{zoufale moc}\nl 
    \opt{lngcs}{akorát}\nl 
    \opt{lngcs}{víc by nezaškodilo}\nl 
    \opt{lngcs}{zoufale málo}}{4}{}
\question{\opt{lngcs}{Matematika a~fyzika mě:}}{%
    \opt{lngcs}{baví, učitel mě motivuje}\nl 
    \opt{lngcs}{baví, učitel mě nemotivuje}\nl
    \opt{lngcs}{nebaví, i~když mě učitel motivuje}\nl
    \opt{lngcs}{nebaví, učitel mě nemotivuje}}{4}{}
\question{\opt{lngcs}{Účastním se některé z~oborových o\-lym\-piád?}}{%
    \opt{lngcs}{fyzikální}\nl
    \opt{lngcs}{matematické}\nl
    \opt{lngcs}{astronomické}\nl
    \opt{lngcs}{logické}\nl
    \opt{lngcs}{biologické}\nl
    \opt{lngcs}{chemické}\nl
    \opt{lngcs}{jiné (vypište)}}{6}{}
\question{\opt{lngcs}{Znám Fyzikální korespondeční seminář (FYKOS)?}}{%TODO
    \opt{lngcs}{ano, od učitele}\nl
    \opt{lngcs}{ano, z~nástěnky ve škole}\nl
    \opt{lngcs}{ano, z~internetu}\nl
    \opt{lngcs}{ano, od kamarádů}\nl
    \opt{lngcs}{ano, dostal jsem dopis}\nl
    \opt{lngcs}{ano, od jinud (odkud?)}\nl
    \opt{lngcs}{ne}\stepcounter{boxcnt}}{5}{}
\question{\opt{lngcs}{Znám jiné korespondenční semináře?}}{%
    \opt{lngcs}{PraSe}\nl
    \opt{lngcs}{FKS}\nl
    \opt{lngcs}{M\&M}\nl
    \opt{lngcs}{KSP}\nl
    \opt{lngcs}{jiný/další (vypište)}}{4}{}
\question{\opt{lngcs}{Řeším FYKOS?}}{%
    \opt{lngcs}{ano}\quad\goto{RESITEL}\nl
    \opt{lngcs}{ne}}{2}{}
\question{\opt{lngcs}{Úlohy FYKOSu jsem:}}{%
    \opt{lngcs}{četl}\nl
    \opt{lngcs}{nikdy neviděl}\quad\goto{CHCI:RESIT}}{2}{}
% otázky pro znechucené neřešitele
\question{\opt{lngcs}{FYKOS neřeším proto, že:}\label{ZNECH:NERES}}{%
    \opt{lngcs}{je moc těžký}\nl
    \opt{lngcs}{je moc lehký}\nl
    \opt{lngcs}{kamarádi by se mi smáli}\nl
    \opt{lngcs}{nemám čas}\nl
    \opt{lngcs}{nelíbí se mi úlohy (proč?)}\nl
    \opt{lngcs}{jiný důvod (vypište)}}{4}{}
{\setlength{\AboxD}{26pt}%



\question{\opt{lngcs}{Chci začít řešit FYKOS?}\label{CHCI:RESIT}}{%
    \opt{lngcs}{ano (vyplňte návratku vloženou v~anketě a~odevzdejte ji spolu s~anketou, zadání vám přijde poštou)}\nl %TODO nebudou mít zadání ve složce?
    \opt{lngcs}{ne}}{2}{}

}

% nová stránka
{\vfil%
\noindent\begin{pspicture}(40pt,40pt)
\psbarcode{\ifnum\teamN<100 0\fi\ifnum\teamN<10 0\fi\the\teamN}{includecheck}{qrcode}%
\end{pspicture}
\hfill\emph{Pokračování na druhé straně.}\vspace{-35pt}\newpage}


\question{\opt{lngcs}{Bavilo by mě soustředění FYKOSu?}\label{SOUS}}{%
    \opt{lngcs}{ano}\nl
    \opt{lngcs}{ne (proč?)}\nl
    \opt{lngcs}{nevím, co od něj mám čekat}\stepcounter{boxcnt}}{1}{
    \noindent\goto{HODNOC}}



% otazky pro stavající řešitele
\question{\opt{lngcs}{Úlohy FYKOSu se mi zdají: (můžete zatrhnout i~více odpovědí)}\label{RESITEL}}{%
    \opt{lngcs}{moc těžké}\nl
    \opt{lngcs}{tak akorát}\nl
    \opt{lngcs}{moc lehké}\nl
    \opt{lngcs}{zajímavé}\nl
    \opt{lngcs}{o~ničem}\nl
    \opt{lngcs}{pracné}\nl
    \opt{lngcs}{příliš matematické}\nl
    \opt{lngcs}{jiné postřehy (vypište)}}{7}{}

%\subsection{Hodnocení Fyziklání}
\medskip\noindent{\sffamily\slshape\Large\opt{lngcs}{Hodnocení Fyziklání}}\smallskip\par
\question{\opt{lngcs}{Na Fyziklání jsem proto, že:}\opt{lngsk}{Na Náboji Junior som, pretože:}\label{HODNOC}}{%
    \opt{lngcs}{mě baví fyzika}\nl 
    \opt{lngcs}{se chci procvičit v~počítání příkladů}\nl 
    \opt{lngcs}{mám volno ze školy}\nl 
    \opt{lngcs}{jiný důvod (vypište)}}{3}{}
\evaluation{\opt{lngcs}{Ohodnoťte jako ve škole:}}{%
    \opt{lngcs}{srozumitelnost zadání úloh}\nl
    \opt{lngcs}{organizace}\nl
    \opt{lngcs}{termín}\nl
    \opt{lngcs}{obtížnost úloh (1--nejlehčí)}
     }{4}{}
\evaluation{\opt{lngcs}{Ohodnoťte jako ve škole opravující:}}{%
    \opt{lngcs}{opravující 0}\nl
    \opt{lngcs}{opravující 1}\nl
    \opt{lngcs}{opravující 2}\nl
    \opt{lngcs}{opravující 3}\nl
    \opt{lngcs}{opravující 4}\nl
    \opt{lngcs}{opravující 5}\nl
    \opt{lngcs}{opravující 6}\nl
    \opt{lngcs}{opravující 7}
     }{8}{}
\question{\opt{lngcs}{Chci se zúčastnit Fyziklání i~příští rok?}}{%
     \opt{lngcs}{ano}\nl
     \opt{lngcs}{ne, příští rok již budu na VŠ}\nl
     \opt{lngcs}{ne, jiný důvod (vypište)}}{2}{}
{\setlength{\AboxD}{35pt}%vetsi boxy
\question{\opt{lngcs}{Které úlohy se mi nejvíce nebo nejméně líbily?}}{%
    \opt{lngcs}{nejvíce}\nl
    \opt{lngcs}{nejméně}}{0}{}
\question{\opt{lngcs}{Která oblast mi v~úlohách chyběla? Které bylo až příliš?}}{%
    \opt{lngcs}{chybělo}\nl
    \opt{lngcs}{nadbytek}}{0}{}
}
%\question{\opt{lngcs}{Měly by mít základní školy v~MFnáboji vlastní kategorii?}\opt{lngsk}{Mali by mať základné školy v~Náboji Junior vlastnú kategóriu?}}{%
%    \opt{lngcs}{ano}\opt{lngsk}{áno}\nl
%    \opt{lngcs}{ne}\opt{lngsk}{nie}}{2}{}
%\question{\opt{lngcs}{Měly by mít týmy složené ze žáků více škol dovolenu účast?}\opt{lngsk}{Mali by mať tímy zložené zo žiakov viacerých škôl povolenú účasť?}}{%
%    \opt{lngcs}{ano}\opt{lngsk}{áno}\nl
%    \opt{lngcs}{ano, ale ve zvláštní kategorii}\opt{lngsk}{áno, ale v~osobitej kategórii}\nl
%    \opt{lngcs}{ne}\opt{lngsk}{nie}}{3}{}

\noindent\begin{minipage}[t]{5mm}\stepcounter{questioNum}\bf\arabic{questioNum}.\end{minipage}
\begin{minipage}[t]{12cm}\sl%
\opt{lngcs}{Podněty, připomínky a~vzkazy organizátorům pište sem.}%
\end{minipage}\nopagebreak\vspace{2mm}%
\par\noindent\fbox{\vrule width \textwidth height 0pt depth 0pt
\vrule width 0pt height 4cm depth 0pt}

\vspace{20pt}
\noindent\begin{pspicture}(40pt,40pt)
\psbarcode{\ifnum\teamN<100 0\fi\ifnum\teamN<10 0\fi\the\teamN}{includecheck}{qrcode}%
\end{pspicture}
\vspace{-70pt}

\bigskip\noindent
\opt{lngcs}{Děkujeme za vyplnění ankety a těšíme se na shledání na příštím ročníku.}%

\hfill\sigstyle{\opt{lngcs}{Organizátoři}\opt{lngsk}{Organizátori}}


\newpage\advance\teamN by 1\repeat%

%\opt{lngcs}{\makefooter}\opt{lngsk}{}%TODO spíše pouze adresa nebo nic



%\noindent{\bf !TODO!} Logičtější uspořádání otázek\par
%\noindent{\bf !TODO!} Logičtější uspořádání odpovědí\par
%\noindent{\bf !TODO!} Větší boxy\par
%\noindent{\bf !TODO!} interpunkce\par
%\noindent{\bf !TODO!} tykat/vykat\par%vykat
%\noindent{\bf !TODO!} nepřehýbat\par
%\noindent{\bf !TODO!} poděkování\par %je tam
%\noindent{\bf !TODO!} vyřešit stránkování\par %snad jo
%\noindent{\bf !TODO!} ??čárový kód\par %beta
%\noindent{\bf !TODO!} ??identifikace org. místa\par% jméno+kód
%\noindent{\bf !TODO!} návod viz MFF\par%hotovo
%\noindent{\bf !TODO!} ??jazykové mutace: CZ/SK\par %připraveno
%\noindent{\bf !TODO!} nesázet v~patičce email pro řešení\par% FIXED v makrech, udělal jsem předtím velkou chybu //Lukas
%\noindent{\bf !TODO!} Vyřešit, kde jsem seděl\par %hotovo
%\noindent{\bf !TODO!} Vyřešit návratku\par%snad jo
%\noindent{\bf !TODO!} ??Návratku do ankety\par %bude vedle
\end{document}
